\cleardoublepage
{
    \chapternonum{附录}

    \appendixsecmajornumbering

    \section{SMART 字段筛选剔除明细}

    表~\ref{tab:app-field-filter-removed} 列出正文~\ref{sec:chap3-field-filter} 节三步字段筛选中剔除的 SMART 读数列、对应剔除原因及适用型号；读数列记法与正文一致（$r_i$ 为原始读数，$n_i$ 为归一化读数）。

    \begin{table}[htbp]
        \centering
        \footnotesize
        \setlength{\tabcolsep}{4pt}
        \caption{\label{tab:app-field-filter-removed}SMART 字段筛选剔除明细}
        \begin{tabularx}{\linewidth}{@{}>{\centering\arraybackslash}p{1.4cm} >{\centering\arraybackslash}p{2.0cm} >{\raggedright\arraybackslash}X >{\centering\arraybackslash}p{0.9cm} >{\centering\arraybackslash}p{0.9cm}@{}}
            \toprule
            步骤 & 读数列 & 剔除原因 & MB1 & MB2 \\
            \midrule
            \multicolumn{5}{@{}l@{}}{\textbf{第一步：恒为空列剔除}} \\
            \addlinespace[2pt]
            第一步 & $r_2$/$n_2$、$r_3$/$n_3$、$r_4$/$n_4$、$r_6$/$n_6$、$r_7$/$n_7$、$r_8$/$n_8$ &
            全体「盘--日」样本均为缺失，属性在公开日志中从未上报 &
            $\checkmark$ & $\checkmark$ \\
            第一步 & $r_{10}$/$n_{10}$、$r_{11}$/$n_{11}$、$r_{13}$/$n_{13}$、$r_{189}$/$n_{189}$、$r_{191}$/$n_{191}$、$r_{193}$/$n_{193}$ &
            同上 & $\checkmark$ & $\checkmark$ \\
            第一步 & $r_{200}$/$n_{200}$、$r_{204}$/$n_{204}$、$r_{205}$/$n_{205}$、$r_{207}$/$n_{207}$、$r_{240}$/$n_{240}$ &
            同上 & $\checkmark$ & $\checkmark$ \\
            \midrule
            \multicolumn{5}{@{}l@{}}{\textbf{第二步：零方差列剔除}} \\
            \addlinespace[2pt]
            第二步 & $r_1$/$n_1$、$r_{170}$/$n_{170}$、$r_{171}$/$n_{171}$、$r_{172}$/$n_{172}$、$r_{173}$/$n_{173}$ &
            存在非空取值但全样本为同一常数，无区分信息 &
            $\checkmark$ & $\checkmark$ \\
            第二步 & $r_{174}$/$n_{174}$、$r_{175}$/$n_{175}$、$r_{188}$/$n_{188}$、$r_{192}$/$n_{192}$、$r_{196}$/$n_{196}$、$r_{197}$/$n_{197}$ &
            同上 & $\checkmark$ & $\checkmark$ \\
            第二步 & $r_{198}$/$n_{198}$、$r_{206}$/$n_{206}$、$r_{211}$/$n_{211}$、$r_{232}$/$n_{232}$、$r_{233}$/$n_{233}$、$r_{245}$/$n_{245}$ &
            同上 & $\checkmark$ & $\checkmark$ \\
            第二步 & $r_{184}$/$n_{184}$、$r_{194}$/$n_{194}$ &
            同上 & --- & $\checkmark$ \\
            第二步 & $r_{244}$/$n_{244}$ &
            同上 & $\checkmark$ & --- \\
            第二步 & $n_{244}$ &
            同上 & --- & $\checkmark$ \\
            \midrule
            \multicolumn{5}{@{}l@{}}{\textbf{第三步：$r$--$n$ 冗余列剔除（仅 $n_i$）}} \\
            \addlinespace[2pt]
            第三步 & $n_5$、$n_9$、$n_{177}$、$n_{181}$、$n_{182}$、$n_{183}$、$n_{187}$、$n_{190}$、$n_{195}$、$n_{199}$ &
            $|\rho_s|>0.8$，$n_i$ 可由 $r_i$ 单调导出 & $\checkmark$ & $\checkmark$ \\
            第三步 & $n_{184}$、$n_{194}$ &
            同上 & $\checkmark$ & --- \\
            \bottomrule
        \end{tabularx}

        \vspace{4pt}
        \begin{minipage}{\linewidth}
            \footnotesize 注：$\checkmark$ 表示该型号执行剔除；「---」表示不剔除。
            第三步 $\rho_s$ 取值见表~\ref{tab:app-raw-norm-spearman}；MB2 上 $r_{184}$、$r_{194}$ 已于第二步因零方差剔除，故第三步不再涉及 $n_{184}$、$n_{194}$。
        \end{minipage}
    \end{table}

    \section{SMART 属性原始读数与归一化读数的 Spearman 相关}

    表~\ref{tab:app-raw-norm-spearman} 列出正文公式~(\ref{equ:spearman-rn}) 所依据的各属性 Spearman 系数：在剔除零方差原始读数列后的有效属性上，于 MB1/MB2 全量日日志按「盘--日」行级配对计算。

    \begin{table}[htbp]
        \centering
        \footnotesize
        \caption{\label{tab:app-raw-norm-spearman}字段筛选后 SMART 原始读数与归一化读数的 Spearman 相关（MB1/MB2）}
        \begin{tabularx}{\linewidth}{c X >{\centering\arraybackslash}X >{\centering\arraybackslash}X >{\centering\arraybackslash}X}
            \toprule
            ID & SMART 属性 & MB1 $\rho_s$ & MB2 $\rho_s$ &  保留列 \\
            \midrule
            5  & 重映射扇区计数           & $-1.0000$ & $-1.0000$ & $r_{5}$ \\
            9  & 通电时间累计             & $-0.9682$ & $-0.9681$ & $r_{9}$ \\
            12 & 通电周期计数             & $-0.0183$ & $-0.0472$ & $r_{12}$/$n_{12}$ \\
            177 & 磨损范围对比值          & $-0.9274$ & $-0.8269$ & $r_{177}$ \\
            180 & 未用的备用块计数        & $0.0407$  & $0.1512$  & $r_{180}$/$n_{180}$ \\
            181 & 编程失败计数            & $-1.0000$ & $-1.0000$ & $r_{181}$ \\
            182 & 擦写失败计数            & $-1.0000$ & $-1.0000$ & $r_{182}$ \\
            183 & 串口降速错误计数        & $-1.0000$ & $-1.0000$ & $r_{183}$ \\
            184 & 点到点错误检测计数      & $-1.0000$ & ---       & $r_{184}$ \\
            187 & 无法校正的错误          & $-1.0000$ & $-1.0000$ & $r_{187}$ \\
            190 & 气流温度                & $-1.0000$ & $-1.0000$ & $r_{190}$ \\
            194 & 温度                    & $-1.0000$ & ---       & $r_{194}$ \\
            195 & 编程错误块计数          & $-1.0000$ & $-1.0000$ & $r_{195}$ \\
            199 & Ultra ATA CRC 错误率    & $-0.9987$ & $-0.9993$ & $r_{199}$ \\
            241 & LBA 写入总数            & $-0.0842$ & $-0.1709$ & $r_{241}$/$n_{241}$ \\
            242 & LBA 读取总数            & $-0.0881$ & $-0.0922$ & $r_{242}$/$n_{242}$ \\
            \bottomrule
        \end{tabularx}

        \vspace{4pt}
        \begin{minipage}{\linewidth}
            \footnotesize 注：$\rho_s$ 按正文公式~(\ref{equ:spearman-rn}) 在同行 $r_i$ 与 $n_i$ 上配对计算；「---」表示该型号上 $r_i$ 与 $n_i$ 至少一侧为零方差而未列入相关统计。正文以 $|\rho_s|>0.8$ 作为冗余判定阈值：超过阈值时仅保留 $r_i$，否则 $r_i$ 与 $n_i$ 均保留。
        \end{minipage}
    \end{table}

    \removeappendixsecmajornumbering
}
